\documentclass[12pt,a4paper]{report}
\usepackage[utf8]{inputenc}
\usepackage[T1]{fontenc}
\usepackage[french]{babel}
\usepackage[french]{nomencl}
\usepackage{lmodern}
\usepackage{graphicx} %Pour les images 
\usepackage{amsmath}
\usepackage{amsfonts}
\usepackage{amssymb}
\usepackage{makeidx}
\usepackage[left=2cm,right=2cm,top=2.5cm,bottom=2.5cm]{geometry}

%Interligne 1.5
\renewcommand{\baselinestretch}{1.5}

%Pour la page de garde
\newcommand{\hsp}{\hspace{20pt}}
\newcommand{\HRule}{\rule{\linewidth}{0.5mm}}

%HyperRef Conf
\usepackage{hyperref}
\hypersetup{
pdftitle={Mémoire de fin d'étude},
colorlinks=true, %colorise les liens
breaklinks=true, %permet le retour à la ligne dans les liens trop longs
urlcolor=black, %couleur des hyperliens
linkcolor=black, %couleur des liens internes
citecolor=black,    %couleur des liens de citations
bookmarksopen=true,
pdftoolbar=false,
pdfmenubar=false,
}

%Gestion des pied et en-tête de page
\usepackage{fancyhdr}
\pagestyle{fancy}
\lhead{}
\chead{}
\rhead{\leftmark}
\lfoot{Alexis BATTAGLI}
\cfoot{Page -\thepage-}
\rfoot{Mémoire de fin d'étude}
\renewcommand{\headrulewidth}{0.4pt}
\renewcommand{\footrulewidth}{0.4pt}

%Gestion des titre et indentation
\usepackage{titlesec}
\renewcommand{\thesection}{\arabic{section}}
\setcounter{secnumdepth}{4} % On affiche une numérotation sur une profondeur de 3
\setcounter{tocdepth}{4}        % La table des matières va a une profondeur de 3
% Alignement des titres :
\titlespacing{\chapter} {0pt} {*0} {*0} {} 
\titlespacing{\section} {4ex} {*0} {*0} {} 
\titlespacing{\subsection} {10ex} {*0} {*0} {} 
\titlespacing{\subsubsection} {18ex} {*0} {*0} {} 

%Gestion des Acronymes & Glossaire
\usepackage[acronym]{glossaries}
\makenoidxglossaries
\loadglsentries{MyGlossaries.tex}
\loadglsentries{MyAcronymes.tex}

\begin{document}

%Page de garde
\begin{titlepage}
  \begin{center}

    \textsc{\LARGE Mémoire de fin d'étude}\\[2cm]

    \textsc{\Large Ingénieur Informatique\\ spécialité Systèmes et Réseaux}\\[1.5cm]

    % Title
    \HRule \\[0.4cm]
    { \huge Conception et réalisation d'un outil de validation d'équipements CWMP\\[0.4cm] }

    \HRule \\[2cm]
    \includegraphics[scale=0.2]{./img/imt_mines_ales-bleu.jpg}
    \includegraphics[scale=0.1]{./img/orange.jpg}
    \\[2cm]

    % Author and supervisor
    \begin{minipage}{0.5\textwidth}
      \begin{flushleft} \large
        \emph{Alternant :} Alexis \textsc{BATTAGLI}\\
        \emph{Maitre d'apprentissage :}Marc \textsc{DOUET}\\
        \emph{Tuteur académique : } Yan \textsc{MORET}
      \end{flushleft}
    \end{minipage}
    \begin{minipage}{0.4\textwidth}
      \begin{flushright} \large
      	\emph{École :} IMT Mines Alès\\
       	\emph{Entreprise :} Orange\\
        \emph{Promotion :} INFRES 7\\
      \end{flushright}
    \end{minipage}

    \vfill

    % Bottom of the page
    {\large Septembre 2014 — Septembre 2017}

  \end{center}
\end{titlepage}
\newpage

\section*{Remerciments}
\newpage
\tableofcontents
\printnoidxglossaries
\listoffigures
\newpage

\section{Introduction}
sklgjegj \gls{tink} zukvhz \gls{tink}.
\subsection{L'entreprise}
\subsection{Le contexte}
\subsubsection{Le Device Mangement à Orange}
\subsubsection{La norme TR-069}
\subsection{Objectifs envisagés}
\subsubsection{Première année}
\subsubsection{Deuxième année}
\subsubsection{Troisième année}

\section{Monté en compétence sur le protocole CWMP}
\subsection{Création d'un ACS Servlet}
\subsection{Études d'équipements}
\subsubsection{Présentation du réseau isolé}
\subsubsection{Test DNS}
\subsubsection{Test de comportement TR-069 d'équipement}
\subsection{Étude de client CWMP}
\subsubsection{Client EasyCWMP}
\subsubsection{Client tr69agent d'Orange}
\subsubsection{Résultats} %conséquence : création du toolkit !
\subsection{Impact sur mon parcours}

\section{Projet principal: Conception et développement d'un outil de test}
\subsection{Contexte}
\subsection{Présentation}
\subsection{Méthode de projet}
\subsection{Travail de préparation}
\subsubsection{Recherche de solution technique}
\subsubsection{Analyse de faisabilité}
\subsection{Conception}
\subsection{Réalisation}
\subsubsection{Travail en équipe}
\subsubsection{Développement}
\subsection{Déploiement}
\subsubsection{Environement}
\subsection{Communication et utilisateur}
\subsection{Livrable du projet}
\subsection{Difficultés, solutions et compétences acquises}
\subsection{Bilan et apport personnel du projet}

\section{Transfert de compétences}

\section{Bilan de compétences}

\section{Conclusion}
\subsection{Atteintes des objectifs}
\subsection{Progression}
\subsection{Synthèse de parcours}
sdgd

\end{document}
